%! Tex main = MarquezSalazarBrandon-PropuestaDeTesis.tex
\subsection{Impacto científico del proyecto}

La contribución fundamental del proyecto radica en la
exploración sistemática del paisaje de vacíos en teoría de
cuerdas tipo IIB mediante métodos inspirados en la naturaleza
(NIM). Al optimizar una función potencial derivada de la
K-teoría, es posible comprender mejor el comportamiento del
espacio de soluciones en compactificaciones con torsión, así
como las condiciones bajo las cuales emergen configuraciones
estables. Esta aproximación computacional permite contrastar,
a nivel numérico, predicciones derivadas del programa
Swampland, tales como la (in)existencia de vacíos de Sitter
estables, la separación de escalas en espacios AdS, o las
condiciones para la presencia de taquiones en el espectro.

Desde el punto de vista metodológico, la comparación del
desempeño de distintos NIM (clasificados según su naturaleza
como algoritmos de búsqueda única o múltiple) sobre un
problema concreto de la física teórica aporta criterios
objetivos para la selección de estrategias de optimización en
futuras investigaciones. En particular, se parte del hecho de
que no existe un método universalmente superior; su eficacia
relativa depende de las características del problema abordado.
La función potencial empleada presenta rasgos (alta
dimensionalidad, múltiples óptimos locales, no linealidad) que
son comunes a otros modelos de la física contemporánea. Así,
la experiencia adquirida en este estudio permite una toma de
decisiones más informada en problemas que requieran
exploración amplia o explotación de regiones prometedoras,
según la naturaleza del espacio de búsqueda.

El software científico derivado del proyecto constituirá una
herramienta para futuras investigaciones, permitiendo
extender, modificar y reproducir los procesos experimentales
aquí desarrollados. Su disponibilidad como código abierto
posibilita que otros investigadores se apoyen en esta base
para abordar problemas similares, ya sea en el mismo dominio
físico o en áreas afines. De este modo, el proyecto participa
del intercambio propio de la comunidad científica a través de
la utilidad concreta de sus productos, que quedan expuestos a
la replicación y el escrutinio.
